\documentclass[11pt]{article}


\setlength{\parindent}{0pt}
\usepackage{xltxtra}
\usepackage{hyperref}
\hypersetup{hidelinks}
\usepackage{url}
\urlstyle{tt}
\usepackage{xcolor}
\definecolor{CVBlue}{RGB}{23,110,191}
\usepackage{calc}
\usepackage{graphicx}
\usepackage{tikz}
\usetikzlibrary{calc}
\usepackage{fontspec}
\usepackage{xeCJK}
\usepackage{enumitem}
\CJKsetecglue{} %% 取消中文与数字之间的间隙


%% 主文档字体设置
\setmainfont[
    Path = fonts/Main/,
    Extension = .otf,
    BoldFont = texgyretermes-bold.otf, % 加粗字体
]{texgyretermes-regular.otf} % 正文字体

% 中文字体设置
\setCJKmainfont[
    Path = fonts/hansans/,
    Extension = .ttf,
    BoldFont = NotoSansSC-Bold.ttf, % 加粗字体
]{NotoSansSC-Regular.otf} % 正文字体


\usepackage{relsize}
\usepackage{xspace}

% 使用 fontawesome（部分图标）
\usepackage{fontawesome} 

% A4纸，上下左右边距
\usepackage[
    a4paper,
    left=1.2cm,
    right=1.2cm,
    top=1.5cm,
    bottom=1cm,
    nohead
]{geometry}

\renewcommand{\baselinestretch}{1.5} % 行间距设为1.5

\usepackage{titlesec}
\usepackage{enumitem}
\setlist{noitemsep} % 取消列表项间的额外间距
%\setlist{nosep} % 取消所有垂直间距
\setlist[itemize]{topsep=0.25em, leftmargin=*}
\setlist[enumerate]{topsep=0.25em, leftmargin=*}

% --- 用于控制【不同项目之间】的垂直距离 ---
\newlength{\interProjectSpacing}
\setlength{\interProjectSpacing}{0.9em} % <--- 在此调整项目之间的距离
\newcommand{\projectsep}{\vspace{\interProjectSpacing}}

% --- 用于控制【项目标题】与下方【项目描述】的距离 ---
\newlength{\intraProjectTitleSep}
\setlength{\intraProjectTitleSep}{0.4em} % <--- 在此调整标题和描述的距离
\newcommand{\titlebreak}{\\[\intraProjectTitleSep]}

% --- 用于控制【项目描述】与下方【要点列表】的距离 ---
\newlength{\intraProjectListTopSep}
\setlength{\intraProjectListTopSep}{0.2em} % <--- 在此调整描述和列表的距离

% =======================================================================


\titleformat{\section}         % 定制 \section 命令 
{\large\bfseries\raggedright} % 将 section 标题设置为大号、粗体且左对齐
{}{0em}                      % 可用于为所有 section 添加前缀（如“章节...”）
{}                           % 可用于在标题前插入代码
[{\color{CVBlue}\titlerule}]  % 在标题后插入一条横线
\titlespacing*{\section}{0cm}{*1.6}{*1.2}



\begin{document}
\pagenumbering{gobble}

%%%% 利用tikz来定位照片
\begin{tikzpicture}[remember picture, overlay] 
    \node[anchor = north east] at ($(current page.north east)+(-2cm,-1.2cm)$) {\includegraphics[height=3cm]{avatar.jpg}};
  \end{tikzpicture}%
  %%%% 利用tikz来定位学校Logo，这里只在第一页显示
  \begin{tikzpicture}[remember picture, overlay] 
    \node[anchor = north west] at ($(current page.north west)+(0.5cm,+1.0cm)$) {\includegraphics[height=6cm]{zju.png}};
  \end{tikzpicture}%
\centerline{\LARGE\bfseries{孙学理}} 

\centerline{\normalsize{\faPhone\ 136-8574-6341 \quad \faEnvelopeO\ \href{mailto:3230100000@zju.edu.cn}{3230100687@zju.edu.cn}}} 

% \centerline{\normalsize{\faGithubSquare\ \href{https://github.com/maksymilan}{https://github.com/maksymilan} \quad \faRssSquare\ \href{https://maksymilan.github.io/}{https://maksymilan.github.io}}} 
    
\section{\makebox[\widthof{\faGraduationCap}][c]{\color{CVBlue}\faGraduationCap}\ 教育背景}    
\textbf{浙江大学} \hfill 2023.9 -- 至今\\[0.5em] % 标题和正文间加一点距离
化学工程与工艺 \quad 大三 
\begin{itemize}[nosep]
    \item 相关课程：《课程综合实践（HPC101）》、《数据结构基础》、《Python》、《C》
    \item 均绩4.2 预计能拿到院内的直博名额
    \item 曾获浙江省政府奖学金，浙江大学二等奖学金、三等奖学金
\end{itemize}

\section{\makebox[\widthof{\faUsers}][c]{\color{CVBlue}\faUsers}\ 项目经历}

% --- 第一个项目 ---
% 将标题行末尾的 \\ 替换为 \titlebreak 命令
\textbf{基于人工智能的化学反应速率预测及其自动化} \hfill 2024.11 -- 至今 \titlebreak
项目描述：旨在构建一辆可自动化运行的小车，在摄像头检测到指定数字或指定溶液颜色时停下。
% 在 itemize 的选项中，使用 topsep=\intraProjectListTopSep 来控制上边距
\begin{itemize}[nosep, topsep=\intraProjectListTopSep]
    \item \textbf{软件}：训练\textbf{CNN模型}进行Mnist数字识别；利用\textbf{贝叶斯优化}进行化学反应速率的预测；使用\textbf{opencv}对溶液变色以及数字黑框进行识别。
    \item \textbf{硬件}：使用\textbf{pymodbus}库进行工业设备通信，控制电机运转和小车运行。
    \item \textbf{核心产出}：，设计了一整套小车运行的方案，并构建了一定的\textbf{可视化操作界面}。
\end{itemize}

% 使用 \projectsep 命令来分隔两个项目
\projectsep

% --- 第二个项目 ---

\textbf{实验流体精密控制系统} \hfill 2025.07 -- 2025.08 \titlebreak
项目描述：一个为化学或材料科学实验设计的自动化流体精密控制系统。它提供了一个图形用户界面（GUI），用于实时监控和控制多个蠕动泵、柱塞泵以及多通道直流电源，并支持自动化实验流程协议。
\begin{itemize}[nosep, topsep=\intraProjectListTopSep]
    \item \textbf{系统架构}：设计并实现了基于\textbf{生产者-消费者模型}的\textbf{多进程（Multi-processing）}架构。通过 multiprocessing.Queue 进行跨进程指令传输与状态同步，有效解决了在处理高频串口 I/O 阻塞时的界面卡顿问题
    \item \textbf{面向对象封装}： 采用\textbf{工厂模式 (Factory Pattern)} 和\textbf{抽象基类 (Abstract Base Class)}构建硬件驱动层。定义了统一的设备接口规范，实现了对不同品牌设备的统一调用，极大地提升了系统的可扩展性与代码复用率。
    \item \textbf{自动化引擎开发}： 独立研发了一套可编排的自动化协议系统。支持用户在 GUI 中\textbf{自定义}复杂的实验时序（启动、参数设置、延时、停止），并基于 \textbf{JSON} 实现了协议序列化存储与加载，实现了实验流程的全自动化执行。
    \item \textbf{交互与可视化}： 基于 \textbf{PyQt6} 构建响应式图形界面，集成 \textbf{PyQtGraph} 实现多通道数据（电压、电流、流速）的毫秒级实时动态绘制；
    利用 \textbf{Pandas} 实现实验数据的后台自动清洗与 Excel 导出。
\end{itemize}

% --- 第三个项目 ---

\textbf{基于历史工业警报数据的分类模型与预测模型} \hfill 2025.09 -- 2025.10 \titlebreak
项目描述：基于工业环境中的警报事件数据集，训练模型预测未来每小时的警报数量。
\begin{itemize}[nosep, topsep=\intraProjectListTopSep]
    \item \textbf{预测模型}：负责\textbf{随机森林}和\textbf{XGBoost}机器学习部分,实现了0.1级别的均方根误差。
    \item \textbf{特征工程}：基于静态时间信息，规划了\textbf{周期性时间特征、滞后性时间特征、短期趋势特征}等高级时间特征，使得模型预测误差降低约99\%。
    \item \textbf{特征筛选}：利用模型自身训练后输出的\textbf{特征重要性（Feature Importance）}，对特征的贡献度进行排序，选取\textbf{平均收益最高}的特征进行模型训练。提高模型训练速度的同时降低过拟合风险。
\end{itemize}

\section{\makebox[\widthof{\faCogs}][c]{\color{CVBlue}\faCogs}\ 技术栈}
\begin{itemize}[nosep]
    \item \textbf{编程语言：} Python, C, 了解过C++和JAVA的部分语法
    \item \textbf{开发工具：} SSH, Git, MakeFile, LaTex
    \item \textbf{操作系统：} Linux（CLI）
\end{itemize}
    
\section{\makebox[\widthof{\faInfo}][c]{\color{CVBlue}\faInfo}\ 其他}
\begin{itemize}[parsep=0.5ex]
    \item \textbf{GitHub：} \href{https://github.com/Shirley-Sun-222}{https://github.com/Shirley-Sun-222} 
\end{itemize}
\end{document}
